% !TEX program = pdflatex
% Minimal article template -- BibLaTeX + Biber.
% Build:  latexmk main.tex        (uses the bundled latexmkrc)
% Clean:  latexmk -c
\documentclass[11pt,a4paper]{article}

% --- Encoding & fonts (pdfLaTeX) --------------------------------------------
\usepackage[T1]{fontenc}
\usepackage[utf8]{inputenc}

% --- Core packages ----------------------------------------------------------
\usepackage[margin=1in]{geometry}
\usepackage{microtype}
\usepackage{amsmath}
\usepackage{amssymb}
\usepackage{graphicx}
\usepackage{booktabs}

% --- Bibliography: BibLaTeX with the Biber backend --------------------------
\usepackage[backend=biber,style=numeric,sorting=nyt]{biblatex}
\addbibresource{refs.bib}

% --- Hyperlinks & smart cross-references ------------------------------------
% Load hyperref late and cleveref last.
\usepackage{hyperref}
\usepackage{cleveref}

% All graphics live under figures/ (relative path).
\graphicspath{{figures/}}

\title{A Minimal Article Template}
\author{Author One\thanks{Affiliation, \texttt{author@example.org}}}
\date{\today}

\begin{document}
\maketitle

\begin{abstract}
This is a placeholder abstract. Replace it with a concise summary of the
work: motivation, method, and the key result.
\end{abstract}

\section{Introduction}
\label{sec:intro}
Placeholder introduction. Cite a source like this~\cite{placeholder-article}.
Cross-reference the equation in \Cref{eq:example}, the table in
\Cref{tab:example}, and the figure in \Cref{fig:example}.

\section{Method}
\label{sec:method}
A single display equation, labelled for reference:
\begin{equation}
  \label{eq:example}
  E = mc^{2}.
\end{equation}

\section{Results}
\label{sec:results}

\begin{table}[htbp]
  \centering
  \caption{A placeholder table using \texttt{booktabs} rules.}
  \label{tab:example}
  \begin{tabular}{lrr}
    \toprule
    Method   & Accuracy & Time (s) \\
    \midrule
    Baseline & 0.81     & 12.4     \\
    Ours     & 0.89     &  9.7     \\
    \bottomrule
  \end{tabular}
\end{table}

\begin{figure}[htbp]
  \centering
  \includegraphics[width=0.6\linewidth]{example.pdf}
  \caption{A placeholder figure. Put \texttt{example.pdf} in \texttt{figures/}.}
  \label{fig:example}
\end{figure}

\section{Conclusion}
\label{sec:conclusion}
Placeholder conclusion.

\printbibliography

\end{document}
